% =========================================================
%  General Formularium – Physics
%  Chapter 6: Electricity & Magnetism
%  Converted from "General Formularium (ENG) - V2.docx"
%  Reviewed and corrected on 21 Oct 2026.
% =========================================================
\documentclass[11pt]{article}
\usepackage[a4paper,margin=2cm]{geometry}
\usepackage{amsmath,amssymb}
\usepackage[dvipsnames]{xcolor}
\setlength{\parindent}{0pt}
\newcommand{\sect}[2]{\section*{\color{#1}#2}}
\newcommand{\sub}[1]{\par\medskip\textbf{#1}\par}
\newcommand{\expl}[1]{{\color{gray}\itshape #1}\par}
\newcommand{\warn}[1]{{\color{BrickRed}\textbf{!} #1}}
\definecolor{cEM}{HTML}{A93226}

\begin{document}

\sect{cEM}{Electricity \& Magnetism}
\expl{Circuits, electric fields, magnetic fields and induction. $U$ = voltage. The unit is given in square brackets.}

% ---------- Circuits ----------
\sub{Electric current}
\[ I = \frac{\Delta q}{\Delta t} \qquad [\mathrm{A}] \]

\sub{Current in a conductor (atomic model)}
\[ I = n\,A\,v\,|q_e| \qquad [\mathrm{A}] \]
\expl{$n$ = number of free electrons per m³, $A$ = cross-section, $v$ = drift velocity.}

\sub{Resistance (Ohm's law) and conductance}
\[ R = \frac{U}{I} \quad [\Omega] \qquad G = \frac{I}{U} = \frac{1}{R} \quad [\mathrm{S}] \]

\sub{Resistance of a wire (resistivity)}
\[ R = \rho\,\frac{l}{A} \qquad R\ [\Omega],\ \rho\ [\Omega\,\mathrm{m}] \]

\sub{Electric power}
\[ P = U I = R I^2 = \frac{U^2}{R} \qquad [\mathrm{W}] \]

\sub{Heat developed in a resistor}
\[ Q = \eta\,U I\,\Delta t = \eta\,\frac{U^2}{R}\,\Delta t = \eta\,I^2 R\,\Delta t \qquad [\mathrm{J}] \]

\sub{Series circuit}
\[ U = U_1 + U_2 + U_3 + \dots \qquad I = I_1 = I_2 = I_3 = \dots \qquad R_s = R_1 + R_2 + \dots \]

\sub{Parallel circuit}
\[ U = U_1 = U_2 = U_3 = \dots \qquad I = I_1 + I_2 + I_3 + \dots \qquad \frac{1}{R_p} = \frac{1}{R_1} + \frac{1}{R_2} + \dots \]

% ---------- Electric field ----------
\sub{Coulomb's law}
\[ F = \frac{1}{4\pi\varepsilon_0}\cdot\frac{|Q_1|\,|Q_2|}{r^2} \qquad [\mathrm{N}] \]

\sub{Electric field strength}
\[ E = \frac{F}{|q|} \qquad E = \frac{1}{4\pi\varepsilon_0}\cdot\frac{|Q|}{r^2} \quad \text{(point charge)} \qquad [\mathrm{N/C}] \]

\sub{Electric potential energy (uniform field)}
\[ E_{\text{pot}} = Q\,E\,d \qquad [\mathrm{J}] \]

\sub{Electric potential and voltage}
\[ V = \frac{E_{\text{pot}}}{Q} \qquad U_{12} = V_1 - V_2 = -\frac{\Delta E_{\text{pot},1\to 2}}{Q} \qquad [\mathrm{V}] \]

\sub{Voltage in a uniform field}
\[ U = E\cdot l \qquad [\mathrm{V}] \]

% ---------- Magnetism ----------
\sub{Magnetic field around a straight conductor}
\[ B = \frac{\mu_0\,I}{2\pi r} \qquad [\mathrm{T}] \]

\sub{Magnetic field inside a solenoid}
\[ B = \frac{\mu_0\,N I}{l} \qquad [\mathrm{T}] \]

\sub{Lorentz force (on a moving charge)}
\[ F_L = |q|\,v\,B\sin\theta \qquad \vec F_L = q\,\vec v\times\vec B \qquad [\mathrm{N}] \]

\sub{Laplace force (on a current-carrying wire)}
\[ F = B\,I\,l\sin\theta \qquad [\mathrm{N}] \]

\sub{Force between two parallel conductors}
\[ F = \frac{\mu_0\,I_1 I_2}{2\pi r}\,l \qquad [\mathrm{N}] \]

\sub{Magnetic flux}
\[ \Phi = B\,A_\perp = B\,A\cos\alpha \qquad [\mathrm{Wb}] \]

\sub{Induced voltage (Faraday's law)}
\[ U_{\text{ind}} = -N\,\frac{d\Phi}{dt} \qquad |U_{\text{ind,avg}}| = N\left|\frac{\Delta\Phi}{\Delta t}\right| = N\left|\frac{\Phi(t_2) - \Phi(t_1)}{t_2 - t_1}\right| \qquad [\mathrm{V}] \]

\end{document}
